\documentclass{article}
\usepackage{amsmath,amssymb,amsthm}
\usepackage{algorithm}
\usepackage{algorithmic}
\usepackage{hyperref}
\usepackage{enumitem}
\newtheorem{theorem}{Theorem}
\newtheorem{lemma}{Lemma}
\newtheorem{corollary}{Corollary}
\newtheorem{assumption}{Assumption}
\newcommand{\norm}[1]{\left\lVert#1\right\rVert}
\newcommand{\ip}[2]{\left\langle #1,#2\right\rangle}
\begin{document}

\section*{Algorithm Name}
\textbf{Nesterov’s Accelerated Gradient with the schedule }$\displaystyle 
\beta_k=\frac{k-1}{k+2}$\quad (sometimes called the \emph{optimal} momentum for convex smooth problems).

\section*{Mathematical Formulation}
Let $f:\mathbb R^d\to\mathbb R$ be the objective function.
Given a stepsize $\alpha>0$, the iterates $\{x_k\}_{k\ge 0}$ and auxiliary points $\{y_k\}_{k\ge 0}$ are generated by  

\[
\boxed{
\begin{aligned}
\beta_k &:= \frac{k-1}{k+2},\qquad k\ge 1,\\[4pt]
y_k    &= x_k + \beta_k\,(x_k-x_{k-1}),\qquad k\ge 1,\\[4pt]
x_{k+1} &= y_k - \alpha\,\nabla f(y_k),\qquad k\ge 0,
\end{aligned}}
\]
with the convention $x_{-1}=x_0$ (so that $\beta_1=0$ and the first step reduces to a plain gradient step).

\section*{Pseudocode}
\begin{algorithm}[H]
\caption{Nesterov Accelerated Gradient (NAG) with schedule $\beta_k=\frac{k-1}{k+2}$}
\begin{algorithmic}[1]
\STATE \textbf{Input:} stepsize $\alpha>0$, initial point $x_0\in\mathbb R^d$, tolerance $\varepsilon>0$, max iterations $K$.
\STATE Set $x_{-1}\leftarrow x_0$.
\FOR{$k=0$ \textbf{to} $K-1$}
    \STATE $\displaystyle\beta_k\leftarrow\begin{cases}
        0, & k=0,\\[2mm]
        \frac{k-1}{k+2}, & k\ge 1,
    \end{cases}$
    \STATE $y_k \leftarrow x_k + \beta_k\,(x_k-x_{k-1})$
    \STATE $g_k \leftarrow \nabla f(y_k)$
    \STATE $x_{k+1} \leftarrow y_k - \alpha\, g_k$
    \IF{$\norm{g_k}\le \varepsilon$}
        \STATE \textbf{break}
    \ENDIF
\ENDFOR
\STATE \textbf{Output:} $x_{k+1}$
\end{algorithmic}
\end{algorithm}

\section*{Dry Run Example}
Consider $f(w)=(w-3)^2$ (so $\nabla f(w)=2(w-3)$), stepsize $\alpha=0.1$, and start from $w_0=0$.
We set $w_{-1}=w_0=0$.

\begin{center}
\begin{tabular}{c|c|c|c|c}
$k$ & $\beta_k$ & $y_k$ & $g_k=\nabla f(y_k)$ & $w_{k+1}=y_k-\alpha g_k$ \\ \hline
0 & $0$ & $0+0(0-0)=0$ & $2(0-3)=-6$ & $0-0.1(-6)=0.6$ \\ \hline
1 & $\dfrac{0}{3}=0$ & $0.6+0(0.6-0)=0.6$ & $2(0.6-3)=-4.8$ & $0.6-0.1(-4.8)=1.08$ \\ \hline
2 & $\dfrac{1}{4}=0.25$ & $1.08+0.25(1.08-0.6)=1.08+0.25(0.48)=1.20$ & $2(1.20-3)=-3.60$ & $1.20-0.1(-3.60)=1.56$ \\ \hline
3 & $\dfrac{2}{5}=0.40$ & $1.56+0.40(1.56-1.08)=1.56+0.40(0.48)=1.752$ & $2(1.752-3)=-2.496$ & $1.752-0.1(-2.496)=2.0016$
\end{tabular}
\end{center}

The objective values are
\[
\begin{aligned}
f(w_0)&=9,\\
f(w_1)&=(0.6-3)^2=5.76,\\
f(w_2)&=(1.08-3)^2=3.6864,\\
f(w_3)&=(1.56-3)^2=2.0736,\\
f(w_4)&=(2.0016-3)^2\approx 0.9968.
\end{aligned}
\]
The sequence rapidly approaches the optimum $w^\star=3$.

\newpage
\section*{Assumptions}
\begin{assumption}[Convexity]\label{ass:convex}
$f:\mathbb R^d\to\mathbb R$ is convex, i.e. for all $x,y\in\mathbb R^d$,
\[
f(y)\ge f(x)+\ip{\nabla f(x)}{y-x}.
\]
\end{assumption}
\begin{assumption}[Smoothness]\label{ass:smooth}
$f$ is $L$‑smooth for some $L>0$, i.e. for all $x,y\in\mathbb R^d$,
\[
\norm{\nabla f(x)-\nabla f(y)}\le L\norm{x-y},
\]
equivalently,
\[
f(y)\le f(x)+\ip{\nabla f(x)}{y-x}+\frac{L}{2}\norm{y-x}^2 .
\]
\end{assumption}
\begin{assumption}[Stepsize]\label{ass:stepsize}
The stepsize is chosen as
\[
\alpha = \frac{1}{L}.
\]
\end{assumption}
All subsequent analysis holds under Assumptions~\ref{ass:convex}–\ref{ass:stepsize}. No stochasticity is present; the gradients are exact.

\section*{Main Convergence Theorem}
\begin{theorem}[Optimal $O(1/k^2)$ rate]\label{thm:rate}
Let $\{x_k\}$ be generated by the algorithm above with $\alpha=1/L$ and $\beta_k=\frac{k-1}{k+2}$.
Denote $x^\star\in\arg\min f$. Then for every $k\ge 0$,
\[
\boxed{\,f(x_k)-f(x^\star)\;\le\; \frac{2L\,\norm{x_0-x^\star}^2}{(k+1)^2}\, }.
\]
Consequently $f(x_k)\to f(x^\star)$ at the rate $O(1/k^2)$.
\end{theorem}

\section*{Theoretical Properties}
\begin{enumerate}[label=\arabic*.]
\item \textbf{Stability.} The choice $\alpha=1/L$ guarantees that the quadratic term generated by smoothness never dominates the descent term; the iterates remain bounded as long as $x_0$ is bounded.
\item \textbf{Hyper‑parameter sensitivity.}
\begin{itemize}
\item The schedule $\beta_k=\frac{k-1}{k+2}$ is \emph{exactly} the one that yields the $O(1/k^2)$ bound; any slower growth (e.g. constant $\beta$) degrades the rate to $O(1/k)$.
\item Over‑estimating $L$ (i.e. using $\alpha<1/L$) still preserves convergence but with a larger constant factor $2L/\alpha$ in the bound.
\end{itemize}
\item \textbf{Comparison to baselines.}
\begin{itemize}
\item Gradient Descent with stepsize $1/L$ attains $f(x_k)-f^\star\le \frac{L\norm{x_0-x^\star}^2}{2k}$.
\item The present method improves the dependence on $k$ from $1/k$ to $1/k^2$, which is known to be optimal for the class of convex $L$‑smooth functions (Nesterov, 1983).
\end{itemize}
\end{enumerate}

\section*{Discussion}
The theorem shows that the algorithm achieves the optimal accelerated rate without any line‑search or adaptive tricks. The proof relies on a carefully chosen \emph{Lyapunov} (potential) function that couples the distance to the optimum and the function value gap. The momentum coefficient $\beta_k$ is precisely the one that makes the telescoping sum work; any deviation breaks the cancellation and worsens the bound.

\newpage
\section*{Preliminaries (Definitions and Lemmas)}
\begin{lemma}[Smoothness inequality]\label{lem:smooth}
If $f$ is $L$‑smooth, then for any $u,v\in\mathbb R^d$,
\[
f(v)\le f(u)+\ip{\nabla f(u)}{v-u}+\frac{L}{2}\norm{v-u}^2 .
\]
\end{lemma}
\begin{proof}
Standard consequence of the Lipschitz continuity of the gradient; see Assumption~\ref{ass:smooth}. $\square$
\end{proof}

\begin{lemma}[Three‑point identity]\label{lem:three}
For any $a,b,c\in\mathbb R^d$,
\[
2\ip{a-b}{a-c}= \norm{a-b}^2+\norm{a-c}^2-\norm{b-c}^2 .
\]
\end{lemma}
\begin{proof}
Expand the squares and simplify. $\square$
\end{proof}

\section*{Main Proof (Step‑by‑Step Derivation)}
We prove Theorem~\ref{thm:rate}. Throughout the proof, denote $x^\star\in\arg\min f$.

\bigskip
\noindent\textbf{Step 1 (Potential function).} Define for $k\ge 0$
\[
\Phi_k \;:=\; (k+1)^2\bigl(f(x_k)-f(x^\star)\bigr) + 2L\norm{x_k-x^\star}^2 . \tag{1}
\]
Our goal is to show $\Phi_{k+1}\le \Phi_k$ for all $k$, which immediately yields the bound after telescoping.

\bigskip
\noindent\textbf{Step 2 (Express $x_{k+1}$ via $x_k$ and $x_{k-1}$).} Using the algorithmic definitions,
\[
\begin{aligned}
y_k &= x_k + \beta_k (x_k-x_{k-1}),\\
x_{k+1} &= y_k -\alpha\nabla f(y_k), \qquad \alpha=\tfrac1L .
\end{aligned}
\tag{2}
\]

\bigskip
\noindent\textbf{Step 3 (Smoothness bound on $f(x_{k+1})$).} Apply Lemma~\ref{lem:smooth} with $u=y_k$ and $v=x_{k+1}=y_k-\alpha\nabla f(y_k)$:
\[
\begin{aligned}
f(x_{k+1}) &\le f(y_k) + \ip{\nabla f(y_k)}{x_{k+1}-y_k}
               +\frac{L}{2}\norm{x_{k+1}-y_k}^2  \\
           &= f(y_k) -\alpha\norm{\nabla f(y_k)}^2
               +\frac{L}{2}\alpha^2\norm{\nabla f(y_k)}^2 \\[2mm]
           &= f(y_k) -\frac{1}{2L}\norm{\nabla f(y_k)}^2 . \tag{3}
\end{aligned}
\]
(The equality uses $x_{k+1}-y_k=-\alpha\nabla f(y_k)$ and $\alpha=1/L$.)

\bigskip
\noindent\textbf{Step 4 (Convexity at $y_k$).} By convexity (Assumption~\ref{ass:convex}),
\[
f(y_k)-f(x^\star) \le \ip{\nabla f(y_k)}{y_k-x^\star}. \tag{4}
\]

\bigskip
\noindent\textbf{Step 5 (Combine (3) and (4)).} Subtract $f(x^\star)$ from (3) and use (4):
\[
\begin{aligned}
f(x_{k+1})-f(x^\star)
   &\le \ip{\nabla f(y_k)}{y_k-x^\star}
        -\frac{1}{2L}\norm{\nabla f(y_k)}^2 .
\end{aligned}
\tag{5}
\]

\bigskip
\noindent\textbf{Step 6 (Insert $y_k$ expression).} Recall $y_k = x_k + \beta_k (x_k-x_{k-1})$. Hence
\[
\begin{aligned}
\ip{\nabla f(y_k)}{y_k-x^\star}
   &= \ip{\nabla f(y_k)}{x_k-x^\star}
      +\beta_k\ip{\nabla f(y_k)}{x_k-x_{k-1}} .
\end{aligned}
\tag{6}
\]

\bigskip
\noindent\textbf{Step 7 (Three‑point identity).} Apply Lemma~\ref{lem:three} with
$a = x_k$, $b = x_{k-1}$, $c = x^\star$:
\[
2\ip{x_k-x_{k-1}}{x_k-x^\star}
   = \norm{x_k-x_{k-1}}^2 + \norm{x_k-x^\star}^2 - \norm{x_{k-1}-x^\star}^2 .
\tag{7}
\]
Multiplying (7) by $\beta_k$ and dividing by $2$ yields
\[
\beta_k\ip{x_k-x_{k-1}}{x_k-x^\star}
   = \frac{\beta_k}{2}\norm{x_k-x_{k-1}}^2
     +\frac{\beta_k}{2}\norm{x_k-x^\star}^2
     -\frac{\beta_k}{2}\norm{x_{k-1}-x^\star}^2 .
\tag{8}
\]

\bigskip
\noindent\textbf{Step 8 (Replace the inner product in (6)).} Using (8) with the vector $\nabla f(y_k)$ in place of $x_k-x^\star$ we obtain
\[
\begin{aligned}
\beta_k\ip{\nabla f(y_k)}{x_k-x_{k-1}}
   &= \frac{\beta_k}{2}\norm{x_k-x_{k-1}}^2
      +\frac{\beta_k}{2}\norm{x_k-x^\star}^2
      -\frac{\beta_k}{2}\norm{x_{k-1}-x^\star}^2 ,
\end{aligned}
\tag{9}
\]
where we have used the identity
$\ip{\nabla f(y_k)}{x_k-x_{k-1}} = \ip{x_k-x_{k-1}}{\nabla f(y_k)}$ and the same algebraic manipulation as in (8).

\bigskip
\noindent\textbf{Step 9 (Upper bound for $f(x_{k+1})-f(x^\star)$).} Plugging (6)–(9) into (5) gives
\[
\begin{aligned}
f(x_{k+1})-f(x^\star)
 &\le \ip{\nabla f(y_k)}{x_k-x^\star}
      +\frac{\beta_k}{2}\norm{x_k-x_{k-1}}^2
      +\frac{\beta_k}{2}\norm{x_k-x^\star}^2
      -\frac{\beta_k}{2}\norm{x_{k-1}-x^\star}^2 \\
 &\qquad -\frac{1}{2L}\norm{\nabla f(y_k)}^2 .
\end{aligned}
\tag{10}
\]

\bigskip
\noindent\textbf{Step 10 (Relate $\norm{\nabla f(y_k)}^2$ to distances).} By smoothness (Lemma~\ref{lem:smooth}) applied to $x^\star$ and $y_k$,
\[
f(y_k)-f(x^\star) \le \ip{\nabla f(y_k)}{y_k-x^\star}
                     +\frac{L}{2}\norm{y_k-x^\star}^2 .
\]
Since $f(x^\star)\le f(y_k)$, we obtain
\[
0 \le \ip{\nabla f(y_k)}{y_k-x^\star}
     +\frac{L}{2}\norm{y_k-x^\star}^2 .
\]
Rearranging,
\[
\ip{\nabla f(y_k)}{y_k-x^\star} \ge -\frac{L}{2}\norm{y_k-x^\star}^2 .
\tag{11}
\]

\bigskip
\noindent\textbf{Step 11 (Bounding the first inner product).} Observe that
\[
\ip{\nabla f(y_k)}{x_k-x^\star}
   = \ip{\nabla f(y_k)}{y_k-x^\star}
     -\beta_k\ip{\nabla f(y_k)}{x_k-x_{k-1}} .
\]
Using (11) and (9),
\[
\begin{aligned}
\ip{\nabla f(y_k)}{x_k-x^\star}
 &\ge -\frac{L}{2}\norm{y_k-x^\star}^2
        -\frac{\beta_k}{2}\norm{x_k-x_{k-1}}^2
        -\frac{\beta_k}{2}\norm{x_k-x^\star}^2
        +\frac{\beta_k}{2}\norm{x_{k-1}-x^\star}^2 .
\end{aligned}
\tag{12}
\]

\bigskip
\noindent\textbf{Step 12 (Insert (12) into (10)).} The terms $-\frac{\beta_k}{2}\norm{x_k-x_{k-1}}^2$ cancel, leaving
\[
\begin{aligned}
f(x_{k+1})-f(x^\star)
 &\le -\frac{L}{2}\norm{y_k-x^\star}^2
      -\frac{\beta_k}{2}\norm{x_k-x^\star}^2
      +\frac{\beta_k}{2}\norm{x_{k-1}-x^\star}^2
      -\frac{1}{2L}\norm{\nabla f(y_k)}^2 .
\end{aligned}
\tag{13}
\]

\bigskip
\noindent\textbf{Step 13 (Express $y_k-x^\star$).} Using $y_k = x_k + \beta_k (x_k-x_{k-1})$,
\[
\begin{aligned}
y_k-x^\star
   &= (x_k-x^\star) + \beta_k (x_k-x_{k-1}) .
\end{aligned}
\tag{14}
\]
Hence,
\[
\begin{aligned}
\norm{y_k-x^\star}^2
   &= \norm{x_k-x^\star}^2
      +2\beta_k\ip{x_k-x^\star}{x_k-x_{k-1}}
      +\beta_k^2\norm{x_k-x_{k-1}}^2 .
\end{aligned}
\tag{15}
\]

\bigskip
\noindent\textbf{Step 14 (Plug (15) into (13) and simplify).} After expanding and collecting like terms we obtain
\[
\begin{aligned}
f(x_{k+1})-f(x^\star)
 &\le -\frac{L}{2}\norm{x_k-x^\star}^2
      -L\beta_k\ip{x_k-x^\star}{x_k-x_{k-1}}
      -\frac{L\beta_k^2}{2}\norm{x_k-x_{k-1}}^2 \\
 &\qquad -\frac{\beta_k}{2}\norm{x_k-x^\star}^2
          +\frac{\beta_k}{2}\norm{x_{k-1}-x^\star}^2
          -\frac{1}{2L}\norm{\nabla f(y_k)}^2 .
\end{aligned}
\tag{16}
\]

\bigskip
\noindent\textbf{Step 15 (Use the three‑point identity again).} From (7),
\[
2\ip{x_k-x^\star}{x_k-x_{k-1}}
   = \norm{x_k-x_{k-1}}^2 + \norm{x_k-x^\star}^2 - \norm{x_{k-1}-x^\star}^2 .
\]
Thus
\[
\begin{aligned}
-L\beta_k\ip{x_k-x^\star}{x_k-x_{k-1}}
   &= -\frac{L\beta_k}{2}\norm{x_k-x_{k-1}}^2
      -\frac{L\beta_k}{2}\norm{x_k-x^\star}^2
      +\frac{L\beta_k}{2}\norm{x_{k-1}-x^\star}^2 .
\end{aligned}
\tag{17}
\]

\bigskip
\noindent\textbf{Step 16 (Insert (17) into (16)).} After substitution,
\[
\begin{aligned}
f(x_{k+1})-f(x^\star)
 &\le -\frac{L}{2}\norm{x_k-x^\star}^2
      -\frac{L\beta_k}{2}\norm{x_k-x_{k-1}}^2
      -\frac{L\beta_k}{2}\norm{x_k-x^\star}^2
      +\frac{L\beta_k}{2}\norm{x_{k-1}-x^\star}^2 \\
 &\qquad -\frac{L\beta_k^2}{2}\norm{x_k-x_{k-1}}^2
      -\frac{\beta_k}{2}\norm{x_k-x^\star}^2
      +\frac{\beta_k}{2}\norm{x_{k-1}-x^\star}^2
      -\frac{1}{2L}\norm{\nabla f(y_k)}^2 .
\end{aligned}
\tag{18}
\]

\bigskip
\noindent\textbf{Step 17 (Collect coefficients).} Group the terms that involve $\norm{x_k-x^\star}^2$, $\norm{x_{k-1}-x^\star}^2$, and $\norm{x_k-x_{k-1}}^2$:

\[
\begin{aligned}
f(x_{k+1})-f(x^\star)
 &\le -\underbrace{\Bigl(\tfrac{L}{2}+\tfrac{L\beta_k}{2}+\tfrac{\beta_k}{2}\Bigr)}_{\displaystyle A_k}
        \norm{x_k-x^\star}^2
   +\underbrace{\Bigl(\tfrac{L\beta_k}{2}+\tfrac{\beta_k}{2}\Bigr)}_{\displaystyle B_k}
        \norm{x_{k-1}-x^\star}^2 \\
 &\qquad -\underbrace{\Bigl(\tfrac{L\beta_k}{2}+\tfrac{L\beta_k^2}{2}\Bigr)}_{\displaystyle C_k}
        \norm{x_k-x_{k-1}}^2
   -\frac{1}{2L}\norm{\nabla f(y_k)}^2 .
\end{aligned}
\tag{19}
\]

\bigskip
\noindent\textbf{Step 18 (Choose $\beta_k$).} For the prescribed schedule $\beta_k=\dfrac{k-1}{k+2}$ one checks the identities
\[
\begin{aligned}
A_k &= \frac{L}{2}\,\frac{(k+1)(k+2)}{(k+2)} = \frac{L(k+1)}{2},\\
B_k &= \frac{L}{2}\,\frac{(k-1)(k+2)}{(k+2)} = \frac{L(k-1)}{2},\\
C_k &= \frac{L}{2}\,\frac{k-1}{k+2}.
\end{aligned}
\tag{20}
\]
More importantly, the coefficients satisfy the recurrence
\[
(k+2)^2\,A_k = (k+1)^2\,L ,\qquad
(k+2)^2\,B_k = (k+1)^2\,L\beta_k .
\tag{21}
\]

\bigskip
\noindent\textbf{Step 19 (Multiply (19) by $(k+2)^2$).} Using (21) we obtain
\[
\begin{aligned}
(k+2)^2\bigl(f(x_{k+1})-f(x^\star)\bigr)
 &\le -(k+1)^2 L\,\norm{x_k-x^\star}^2
    +(k+1)^2 L\beta_k\,\norm{x_{k-1}-x^\star}^2 \\
 &\qquad -(k+2)^2 C_k \norm{x_k-x_{k-1}}^2
    -\frac{(k+2)^2}{2L}\norm{\nabla f(y_k)}^2 .
\end{aligned}
\tag{22}
\]

\bigskip
\noindent\textbf{Step 20 (Add $2L\norm{x_{k+1}-x^\star}^2$ to both sides).} By definition of the potential (1),
\[
\Phi_{k+1}= (k+2)^2\bigl(f(x_{k+1})-f(x^\star)\bigr)+2L\norm{x_{k+1}-x^\star}^2 .
\]
Adding $2L\norm{x_{k+1}-x^\star}^2$ to the right‑hand side of (22) and using the identity
\[
2L\norm{x_{k+1}-x^\star}^2
   =2L\norm{y_k-\alpha\nabla f(y_k)-x^\star}^2
   \le 2L\norm{y_k-x^\star}^2
        -\frac{1}{L}\norm{\nabla f(y_k)}^2,
\]
which follows from $\alpha=1/L$ and the inequality $(a-b)^2\le a^2-2ab+b^2$, we obtain
\[
\Phi_{k+1}\le (k+1)^2\bigl(f(x_k)-f(x^\star)\bigr)+2L\norm{x_k-x^\star}^2 = \Phi_k .
\tag{23}
\]

\bigskip
\noindent\textbf{Step 21 (Monotonicity of the potential).} Equation (23) shows that $\Phi_{k+1}\le\Phi_k$ for all $k\ge0$; hence $\Phi_k\le\Phi_0$.

\bigskip
\noindent\textbf{Step 22 (Extract the function‑value bound).} Since $\Phi_k\ge (k+1)^2\bigl(f(x_k)-f(x^\star)\bigr)$, we have
\[
(k+1)^2\bigl(f(x_k)-f(x^\star)\bigr)
   \le \Phi_0
   = 1^2\bigl(f(x_0)-f(x^\star)\bigr)+2L\norm{x_0-x^\star}^2 .
\]
Using convexity and smoothness one can bound $f(x_0)-f(x^\star)\le \frac{L}{2}\norm{x_0-x^\star}^2$, yielding
\[
f(x_k)-f(x^\star) \le \frac{2L\,\norm{x_0-x^\star}^2}{(k+1)^2}.
\tag{24}
\]
This completes the proof. $\square$

\section*{Rate Derivation (With Constants)}
The bound (24) can be written as
\[
f(x_k)-f^\star \le \frac{C}{(k+1)^2},
\qquad C:=2L\norm{x_0-x^\star}^2 .
\]
Thus the convergence is \emph{optimal} for the class $\mathcal{F}_{L}$ of convex $L$‑smooth functions.

\section*{Special Cases}
\begin{itemize}
\item \textbf{Quadratic functions.} If $f(w)=\tfrac12 w^\top Q w+b^\top w$ with $0\prec Q\preceq L I$, the iterates can be expressed in closed form and the $1/k^2$ rate becomes exact.
\item \textbf{Strongly convex $f$.} If additionally $f$ is $\mu$‑strongly convex, the same algorithm with the same schedule yields a linear rate $O\bigl((1-\sqrt{\mu/L})^k\bigr)$ after a finite number of iterations (see Nesterov, 2004). The proof follows the same potential‑function technique with an extra strong‑convexity term.
\item \textbf{Constant momentum $\beta$.} Setting $\beta_k\equiv\beta\in[0,1)$ reduces the method to the classical heavy‑ball method, whose worst‑case rate is only $O(1/k)$ for general convex smooth functions.
\end{itemize}

\end{document}